\documentclass[../../Orator.tex]{subfiles}

\begin{document}

The attempt to find an equation which can accurately describe the behavior of a chosen system is one of the most fundamental aspects of mathematics.

A model is a series of mathematical equations capable of replicating the behavior of a system. A role of these models, known as dynamic systems, is to bring light in the gaps of understanding and explain the underlying mechanisms behind some function. 
The construction of a model requires equations complex enough to accurately describe the dynamics of interest yet, preferably, simple enough so that mathematical tools exist to analyze the equations\footnotemark. 
\footnotetext{The best model of a cat is a cat. Preferably the same cat\footnotemark.-`\textit{Philosophy of Science, 1945, Arturo Rosenblueth \tbr{1900-1970}}'}
There exists any number of limitations in the pursuit of this task, some inherent to the system, others inherent to our modern construction of mathematics.
\footnotetext{If man could be crossed with the cat, it would improve man, but it would deteriorate the cat.-`\textit{Mark Twain}'}

\section{Dynamic Systems}

When looking at the brain in abstract form it can be seen as a multi-dimensional dynamical system governed by system variables and parameters.

Fundamentally, a dynamical system is a system that evolves over time. This system is dependent on fixed rules which govern the evolution of the system. Thus a dynamical system is the opposite of a static system, which, unlike dynamical systems, does not depend on the past inputs. A simple example of a dynamical system can be that of a pendulum swinging back and forth \cite{}. 

There are different types of dynamical systems. A system can be deterministic or stochastic, discrete or continuous, linear or nonlinear, and autonomous or non-autonomous \cite{}. 

If a system is deterministic, it is possible to predict each following state given an initial state. It is also possible for a system to have an element of randomness to it. In this case, the system is said to be stochastic \cite{}. 

A discrete system is one where there is only measured a position of the integer values of time. On the other hand, is a continuous system, where the positions are measured continuously (ie. for every possible time) \cite{}.

A linear system is one where it is only composed of linear functions and a nonlinear system contains at least one nonlinear component \cite{}.

Lastly, there is the difference between an autonomous system, which is one where it does not depend on the independent variable \(t\) \cite{}, and a non-autonomous system, which does depend on \(t\) \cite{}.

%The 
%Some of these parameters

\begin{comment}
    \begin{split}\left[\begin{array}{ccll}
    {\displaystyle \frac{du}{dt}} &=& u\left(1-u^{2}\right)-w+I \equiv F(u,w)\\[.2cm]
    {\displaystyle \frac{dw}{dt}} &=& \varepsilon \left(u -0.5w+1\right) \equiv \varepsilon G(u,w)\, ,\\
    \end{array}\right.\end{split}
    [ref:https://neuronaldynamics-exercises.readthedocs.io/en/latest/exercises/brunel-network.html]
\end{comment}

\textbf{Phase space} is the space spanned by the system's variables and the trajectories it contain are those variables traversing time

\textit{Bifurcation} designates a qualitative change in the dynamical behavior of the system associated with modifications in the system parameters~\cite{STEFANESCU2012748}.




\subsubsection{Dimensionality reduction}
A neuron model, such as the Hodgkin-Huxley model, uses a system of four differential equations

\subsection*{Networks}
A complex system can be described by a network or a graph with complex topology, whose nodes are the elements of the system and whose edges represent the interactions among them. `One significant recent discovery in the field of complex networks is the observation that a number of large-scale and complex networks are scale-free, that is, their connectivity distributions have the
power-law form'~\cite{wang2002synchronization}.




\subsection*{Taxonomy of seizures}


When modeling neuronal interactions, and in particular modeling of seizures, the model must reflect the fact that thousands of neurons need to interact in order to display seizure-like activity. 
The first major roadblock becomes the quantity of neurons involved, it is not possible to measure the activity of every individual neuron, even assuming one could do so, it would be a monumental task to extract useful information from the bulk. 
Therefore we must be smart about creating a model, choosing the `right' approach to simplify the equations can save astronomical amounts of time and effort.
\begin{enumerate}
    \item \textbf{Model average activity}
    \item \textbf{Reduce parameter space}
    \item \textbf{Tackle a smaller problem}
\end{enumerate}


A system can be modeled deterministicly or contain stochasticly. When collecting data, measurements will rarely be completely deterministic. This is not because forces act in unpredictable manner, rather that there exists far too many elements in the `real' world for any model to fully account for them. 
As a consequence, many models will incorporate both deterministic and stochastic elements to more accurately reflect the observed data.
\begin{enumerate}
    \item \textbf{Model simplification}
    \item \textbf{Stochastic activity}
    \item \textbf{Stochastic Inputs}
\end{enumerate}



\newpage
\section{Physiological Parameters}
It's important to determine what is even being modeled, which features can be parameterized and which can not.


\newpage
\section{Building up to Hodgkin-Huxley}

The Hodgkin-Huxley Model of Neuron Action Potential is regarded as being a part of the great achievements of 20th-century biophysics. Receiving the Nobel Prize in Physiology or Medicine for what is inarguably an incredible feat of human ingenuity. However, the authors of this paper are students, and therefore such a sophisticated model might be beyond the direct reach of abilities within the given timeframe. 
As such, the exploration begins by investigating simplified Hodgkin-Huxley derived models to see what insights can be built.

\subsection{FitzHugh-Nagumo}
`The models in this category are highly simplified toy models that qualitatively describe the membrane voltage as a function of input. They are mainly used for didactic reasons in teaching but are not considered valid neuron models for large-scale simulations or data fitting'.\footnote{Wikipedia} 
%\subsubsection{Didactic toy model}
\begin{subequations}\label{eq:FitNagSys}
    \begin{align}
        \ode{V}{t} &= V - \frac{V^3}{3} - w + I_{ext} \\ 
        \tau\ode{w}{t} &= V - a - b\,w 
    \end{align}
\end{subequations}


\subsection{Chay}
In 1985, T.R. Chay proposed a model of three-dimensional nonlinear differential equations based on the HH model to study chaotic behavior and show ionic events in excitable membranes. 
\begin{subequations}\label{eq:ChaySys}
    \begin{align}
        \ode{V}{t} &= g_\rmm{I}  m^3_\infty h_\infty \br{V_\rmm{I} - V} + g_\rmm{K, V} n^4 \br{V_\rmm{K} - V} + g_\rmm{K, C}  \frac{C}{1+C}\br{V_\rmm{K} - V} + g_\rmm{L} \br{V_m - V_\rmm{L}} \\ 
        \ode{n}{t} &= \frac{n_\infty - n}{\tau_n} \\
        \ode{C}{t} &= \rho \, \sbr{m^3_\infty h_\infty \br{V_c - V} - k_C C}
    \end{align}
\end{subequations}

where \textit{V}, \textit{n}, and \textit{C} are membrane potential, probability of the voltage-sensitive \(\Kp\) channel, and intracellular concentration of \(\Capp\) ions, respectively. The Chay model parameters are adopted from \textit{`paper'} and collected in Table

The \(m_\infty\), \(h_\infty\), and \(n_\infty\) are calculated by \(y_\infty = \alpha_y / \br{\alpha_y + \beta_y} \) formula, and the explicit expressions for 
\(\alpha_m, \beta_m, \alpha_h, \beta_h, \alpha_n, \beta_n\), and \(\tau_n\) are given by:
\begin{align*}
    \alpha_m &= 0.1 \frac{ 25 + V }{1 - \exp{-0.1 \, V - 2}}, &
    \alpha_h &=  0.07 \exp{-0.05\,V -2.5}, &
    \alpha_n &= 0.01 \, \frac{ 20 + V }{ 1 + \exp{-0.1 \, V - 2}} \\
    \beta_m  &= 4 \exp{-\br{\frac{ V + 50 }{ 18 } } }, &
    \beta_h  &= \frac{ 1 }{ 1 + \exp{-0.1 \, V - 2}}, &
    \beta_n  &= 0.125 \exp{- \frac{V + 30}{80}}, \\
    \tau_n &= \frac{1}{ r_n \, \br{\alpha_n + \beta_n} }
\end{align*}



\newpage
\subsection{Hodgkin-Huxley}\label{sec:HHMeth}

The Hodgkin-Huxely model of Action Potential is a system of non-linear differentiable equations with four state variables with respect to time, \(V_m\br{t}, \ n\br{t}, \ m\br{t}, \ h\br{t}\)~\cite{HodHux1952}. The model is built from approximating the characteristics of excitable cells, such as neurons, to a circuit-like construct \cref{fig:MembraneCircut}. 
Mathematically, the flow of current is represented as:
    \begin{align}
        I_c &= C_m \ode{V_m}{t} \\
        \intertext{and the current through a given ion channel is the product of that channel's conductance and the reversal potential for the specific ion}
        I_\bfm{ion} &= {g_\bfm{ion}}(V_\rmm{m}-V_i)  
    \end{align}
        where \(V_\bfm{ion}\) is the reversal potential of the specific ion channel. Thus, for a cell with sodium and potassium channels, the total current through the membrane is given by:
    \begin{subequations}\label{eq:HodHuxSys}
        \begin{equation}
            I = C_m \ode{V_m}{t} + \bar{g}_\rmm{K} n^4 \br{V_m - V_\rmm{K}} + \bar{g}_\rmm{Na} m^3 h \br{V_m - V_\rmm{Na}}  + \bar{g}_\rmm{L} \br{V_m - V_\rmm{L}} 
        \end{equation}
The duo of Hodgkin and Huxley divined the model from observing the smooth current changes in terms of pores or channels that were either open or closed, and by using a statistical approach generated predictions for the probability of channels being open. 

Accounting for the current through a any given ion channel is the product of the channel's conductance and equilibrium potential for the specific ion.
These `gating' variables \(m\) and \(h\), or \(n\) describe the time dependent kinetics of the voltage
        \begin{align}
            \ode{n}{t} &= \alpha_n \br{V_m} \br{1-n} - \beta_n \br{V_m} n \\
            \ode{m}{t} &= \alpha_m \br{V_m} \br{1-m} - \beta_m \br{V_m} m \\
            \ode{h}{t} &= \alpha_h \br{V_m} \br{1-h} - \beta_h \br{V_m} h 
        \end{align}
    \end{subequations}

The ion channel activation/inactivation\footnotemark~probabilities, denoted by \(\alpha_p, \, \beta_p : \, \br{n,m,h} \in p\), are defined such that:
\begin{align}
    \alpha_p\br{V_m} &= p_\infty \br{V_m} / \tau_p \\
    \beta_p\br{V_m}  &= \br{1 - p_\infty \br{V_m}} / \tau_p 
\end{align}
With \(p_\infty\) and its inverse \(1-p_\infty\) being the steady state values for activation and inactivation respectively~\cite{}. 
\footnotetext{\underline{a}ctivation gives \(\alpha\), \underline{b}nactivation gives \(\beta\), really these are the obvious choices}


In the original paper by Hodgkin and Huxley, the relationships of \(\alpha_p, \text{ and }\, \beta_p\) were defined as:
\begin{align*}
    \alpha_n \br{V_m} &= 0.01 \, \cfrac{ 10 - V }{ \exp{ 10 - V } - 1}, &
    \alpha_m \br{V_m} &= 0.1 \cfrac{ 25 - V }{\exp{\cfrac{25-V}{10}} - 1}, &
    \alpha_h \br{V_m} &=  0.07 \, \exp{-\cfrac{V}{20}}, \\
    \beta_n \br{V_m}  &= 0.125 \, \exp{-\cfrac{V}{80}}, &
    \beta_m \br{V_m}  &= 4 \, \exp{-\cfrac{V}{18}}, &
    \beta_h \br{V_m}  &= { \cfrac{1}{\exp{\cfrac{30 - V}{10}} + 1}},
\end{align*}
Where \(V = V_\rmm{rest} - V_m\) represents the polarization in \unit{\milli\volt}

In order to characterize the ion-channels, the equations can be fitted to voltage clamp\footnotemark~data.

\footnotetext{An assay that measures the flow of current through a neuronal cell membrane by `clamping' the potential at an unchanging value.}

\begin{comment}
    Breaking down the equation we find the following regions:
    \begin{equation*}
        I = 
        \overbrace{C_m \ode{V_m}{t}}^{\mathclap{\text{rate of change of voltage scaled by capacitance}}} + 
        \bar{g}_\rmm{K} n^4 \br{V_m - V_\rmm{K}} + 
        \bar{g}_\rmm{Na} m^3 h \br{V_m - V_\rmm{Na}}  + 
        \bar{g}_\rmm{L} \br{V_m - V_\rmm{L}} 
    \end{equation*}




    \begin{figure}[h]
        \centering
        \begin{tikzpicture}
            
            \draw (0,0) circle (0.25);
            
        \end{tikzpicture}
        \caption{fart}\label{<label>}
    \end{figure}
\end{comment}


\end{document}